\section{背景知识：为什么需要内核堆？}

PMM 以 4\,KiB 页为单位分配，VMM 的 \texttt{vmm\_alloc\_pages} 同样以页为粒度。
但内核中大量数据结构远小于一页——进程控制块 128\,B、VMA 节点 32\,B、命令表项十几字节。
若每个小对象都分配整页，内部碎片率高达 $1 - 128/4096 \approx 96.9\%$。

\begin{whybox}
为什么不直接用全局数组（静态分配）？

静态数组在编译时就固定了容量，无法按运行时需求伸缩；为"以防万一"必须预留最大容量，浪费严重；
每新增一种数据结构就要新增一个数组，代码耦合也越来越重。
内核堆（\texttt{kmalloc}/\texttt{kfree}）就是给内核提供任意大小、动态生命周期的分配能力——
类比用户态的 \texttt{malloc}/\texttt{free}。
\end{whybox}

% ============================================================================

